% 原书分章 PDF 第 2—4 页：章首导言与第 11.1 节。分部页由主文件管理。
\chapter{累积疫苗筛选效应的靶向估计}
\label{ch:cumulative-vaccine-sieve-effects}
\chapterauthors{David Benkeser、Marco Carone 与 Peter Gilbert}
\chapteraffiliations{%
D. Benkeser（通讯作者）：美国埃默里大学罗林斯公共卫生学院生物统计学与生物信息学系，1518 Clifton Rd. NE, Atlanta, GA 30322, USA；电子邮箱：\href{mailto:benkeser@emory.edu}{benkeser@emory.edu}。\\
M. Carone：美国华盛顿大学生物统计学系，F-644 Health Sciences Building, 1705 NE Pacific Street, Seattle, WA 98195, USA；电子邮箱：\href{mailto:mcarone@uw.edu}{mcarone@uw.edu}。\\
P. Gilbert：美国弗雷德·哈钦森癌症研究中心疫苗与传染病部，1100 Fairview Ave, Seattle, WA 98109, USA；电子邮箱：\href{mailto:pgilbert@scharp.org}{pgilbert@scharp.org}。\\
本章 DOI：\url{https://doi.org/10.1007/978-3-319-65304-4_11}。}

过去一个世纪，人们已经研制出有效疫苗，用于预防多种病原体所致的疾病。然而，预防人类免疫缺陷病毒（human immunodeficiency virus，HIV）感染的有效疫苗仍未研制成功。HIV 的遗传特征具有显著异质性，这是研制预防 HIV 感染疫苗所面临的一项挑战。预防性 HIV 疫苗通常只由少数几种抗原构成，对于与疫苗所含抗原相似的病毒株引起的感染，可能具有良好的保护作用；但对于抗原性不同的病毒株所致疾病，却可能无法提供保护。因此，在评价预防性 HIV 疫苗时，需要研究疫苗效力（vaccine efficacy）是否随病毒特征而变化，以及如何变化；这一研究领域称为筛选分析（sieve analysis）（Gilbert 等，1998，2001）。可以把疫苗想象成一个筛子：它诱导针对特定病毒株的免疫力，形成阻挡感染的屏障；同时，也可能存在“筛孔”，即能够突破疫苗屏障的 HIV 病毒株。

\begin{sourcebox}
在某个给定遗传位点上，疫苗筛选效应（vaccine sieve effect）定义为：比较在该位点与疫苗匹配的病毒和不匹配的病毒时，疫苗效力的差异。识别筛选效应，有助于指导未来疫苗所应包含的抗原选择，这些疫苗可能是多价疫苗。多价疫苗能够对遗传上多样的病毒提供更广泛的保护，因而可能具有更高的总体效力。
\end{sourcebox}

从统计角度看，筛选分析在竞争风险（competing risks）框架下进行。在存在竞争风险的情形中，试验参与者面临发生若干相互竞争的终点事件的风险。HIV 疫苗试验根据引起感染的病毒基因型，定义不同的终点；每种基因型对应一种单独的终点类型。在竞争风险情形中，评价处理对给定类型终点风险的影响，通常使用瞬时参数或累积参数。选择哪一种参数，取决于具体的科学背景（Pintilie，2007）；两者都曾用于评价疫苗筛选效应（Gilbert，2000）。瞬时参数通常基于原因特异性风险函数（cause-specific hazard function）（Prentice 等，1978；Benichou 和 Gail，1990；Gaynor 等，1993；Lunn 和 McNeil，1995），它定义为：在尚未发生任何终点事件的人群中，发生给定类型终点事件的瞬时概率。最常研究的累积参数是累积发生率（cumulative incidence），即某事件在给定时间之前发生、且属于特定类型的概率。虽然两种参数都与疫苗筛选效应的评价有关，但当疫苗效应随时间减弱时，累积参数可能具有更大的公共卫生意义。

筛选分析通常使用 Aalen--Johansen 估计量，对累积发生率进行推断（Aalen，1978）。只需很少的假设，这一估计量就能具有若干理想性质。只要删失是非信息性的，它就是一致的。此外，若没有测量任何预后协变量，它还具有渐近非参数有效性。然而，任何纵向研究通常都需要考虑信息性删失的问题；HIV 疫苗试验也会常规收集性行为风险等预后协变量。利用这些协变量，可以减弱对删失机制的假设，并提高效率。

人们已经提出一些半参数方法，将协变量纳入竞争风险数据的分析，并应用于筛选分析。其中包括针对原因特异性风险的比例风险回归（Prentice 等，1978；Lunn 和 McNeil，1995），以及针对亚分布风险的比例风险回归（Fine 和 Gray，1999）。这些基于风险函数的方法，可以利用风险函数与发生率之间的已知关系，计算累积发生率的估计；不过，它们更常用于估计基于风险函数的效力参数。这些半参数方法的一个缺点是，无论估计瞬时参数还是累积参数，都要求正确设定一个有限维回归模型。

\begin{sourcebox}
模型设定错误时，目标参数通常难以解释。例如，已知设定错误的原因特异性 Cox 模型，其估计目标会涉及删失分布（Struthers 和 Kalbfleisch，1986）；设定错误的亚分布风险模型也存在同样的问题（Grambauer 等，2010）。删失分布通常被视为某项研究特有的干扰因素，而不是我们感兴趣的总体特征。因此，估计目标涉及删失分布，在评价处理效力时是一种不理想的性质（Stitelman 和 van der Laan，2011）。
\end{sourcebox}

对于只有一种终点类型的情形，Moore 和 van der Laan（2009a）以及 van der Laan 和 Gruber（2012）提出了边际生存概率的靶向估计量。这些研究提出并评价了利用超级学习器（super learner）等基于机器学习的估计量进行协变量调整的方法。研究表明，这些估计量能够提高效率，并增强对信息性删失的稳健性。本章将说明，如何将 van der Laan 和 Gruber（2012）提出的靶向最小损失估计（targeted minimum loss-based estimation，TMLE）推广到竞争风险情形，在筛选分析中估计累积发生率。我们将利用一项近期开展的预防性 HIV 疫苗效力 II 期试验的数据，说明这些方法。

\section{观测数据}

我们考虑一项预防性 HIV 疫苗效力试验。试验招募 $n$ 名个体，并测量每个人的一组基线特征 $L_0$，这组特征可能具有很高的维数。随后，个体被分配接受试验疫苗 $A_0=1$ 或对照疫苗 $A_0=0$，分配也可能依据 $L_0$ 进行。试验参与者需要按照预定时间，参加 $K+1$ 次定期临床访视，接受 HIV 感染检测。我们考虑某个特定遗传位点受到关注的情形。对 $k=1,\ldots,K+1$，定义二元感染指示向量 $L_k=(L_{k,1},L_{k,2})$：若参与者在第 $k$ 次访视或此前，感染的病毒在目标位点上与疫苗中的参考病毒匹配，则 $L_{k,1}=1$；否则 $L_{k,1}=0$。类似地，若参与者在第 $k$ 次访视或此前，感染的病毒与疫苗中的参考病毒不匹配，则 $L_{k,2}=1$；否则 $L_{k,2}=0$。随访期间，一些参与者可能撤回知情同意，或因其他原因退出研究。对 $k=1,\ldots,K$，用 $A_k$ 表示参与者是否参加第 $k+1$ 次临床访视。若参与者缺席某次临床访视，就将其视为右删失。因此，观测数据可以表示为下列随机变量的 $n$ 个独立副本：
\[
 O=(L_0,A_0,L_1,A_1,\ldots,A_K,L_{K+1})\sim P_0.
\]
我们不对 $P_0$ 作任何假设，因此统计模型 $\mathcal M$ 是非参数模型。

需要说明，对 $k=1,\ldots,K$，也可以允许 $L_k$ 包含参与者随时间变化的特征；不过，我们的数据分析并未纳入这些特征。用 $\bar L_k=(L_0,\ldots,L_k)$ 和 $\bar A_k=(A_0,\ldots,A_k)$，表示时变过程中截至给定时点 $k$ 的历史。还用 $\bar0_m$ 和 $\bar1_m$，分别表示长度为 $m$ 的全零向量和全一向量。
